\section*{ID8759: Relation: Upsilonin+}
\paragraph{Metadata.}
PiID: 0249011419521 | RTEID: RTE:P34013.1422:T5.7777 | UUID: 9b396753-0d77-5ff5-b2d3-dfe8021c4aae

\paragraph{Canonical Statement.}
ace, but each field line still has a branch identity with respect to its TZP. The coordinate system could be extended to label multiple TZPs, but that’s beyond our scope here. We focus on one TZP and its two branches. One way to formalize DAANSsphere is to say it’s the disjoint union of two spheres: \$\textbackslash{}mathcal\{S\} \{DAANS\} = S\textasciicircum{}2 \{(+)\} \textbackslash{}sqcup S\textasciicircum{}2\_\{(-)\}\$ . A point on \$\textbackslash{}mathcal\{S\} \{DAANS\}\$ is \$(\textbackslash{}hat\{n\},\textbackslash{}Upsilon)\$ where \$\textbackslash{}hat\{n\}\$ is a unit direction vector in \$ℝ\textasciicircum{}3\$ (point on \$S\textasciicircum{}2\$ ) and \$\textbackslash{}Upsilon\textbackslash{}in\{+,-\}\$ . We can also embed this into a single mathematical object by introducing a function like \$\textbackslash{}sigma = \textbackslash{}pm 1\$ and consider \$(\textbackslash{}sigma \textbackslash{}hat\{n\})\$ in an augmented space. However, since we likely want to do calculus on this, we can think in terms of covering space: \$\textbackslash{}mathcal\{S\} \{DAANS\}\$ is a 2-sheeted cover of \$S\textasciicircum{}2\$ . Now, how does \$\textbackslash{}mathcal\{S\} \{DAANS\}\$ help in the theory? Primarily, it allows us

\paragraph{Canonical Equation.}
\[
\Upsilon\in{+,-}
\]

\paragraph{Dictionary.}
Relation: Upsilonin+ — Upsilonin+,-

\paragraph{Encyclopedia.}
Relation: Upsilonin+ is a symbolic expression, written Upsilonin+,-. It introduces a symbol or relation used elsewhere in the framework. Within the Trawin Zero Point registry it serves as one element of the framework's mathematical narrative.
